% Generated draft: requires mathematical and source review.
\documentclass[11pt]{article}
\usepackage[margin=25mm]{geometry}
\usepackage{fontspec}
\setmainfont{Latin Modern Roman}
\usepackage{amsmath,amssymb,mathtools,unicode-math}
\setmathfont{Latin Modern Math}
\usepackage{enumitem,needspace,xurl,hyperref}
\setlength{\parindent}{0pt}
\setlength{\parskip}{5pt}
\setlength{\emergencystretch}{3em}
\newcommand{\N}{\mathbb{N}}
\newcommand{\Z}{\mathbb{Z}}
\newcommand{\Q}{\mathbb{Q}}
\newcommand{\R}{\mathbb{R}}
\newcommand{\C}{\mathbb{C}}
\newcommand{\F}{\mathbb{F}}
\newcommand{\A}{\mathbb{A}}
\newcommand{\PP}{\mathbb{P}}
\newcommand{\E}{\mathbb{E}}
\newcommand{\eps}{\varepsilon}
\newcommand{\dd}{\,\mathrm{d}}
\begin{document}
\section*{UnsolvedMath ID: 11000047}
\textbf{Source problem code:} catalogue entry
\subsection{Definitions and mathematical statement}
Let \(\mathfrak H_g\) be Siegel upper half-space and
\[
 \mathcal A_g=\operatorname{Sp}(2g,\Z)\backslash\mathfrak H_g,
 \qquad \mathfrak H_g\cong\operatorname{Sp}(2g,\R)/U(g),
\]
the usual moduli orbifold of principally polarized complex abelian varieties. Let \(\mathcal J_g\subseteq\mathcal A_g\) denote the source's period locus. In the usual smooth-curve interpretation it is the image of the Torelli period map from \(\mathcal M_g\), not automatically its closure. A real symplectic matrix \(S=\left(\begin{smallmatrix}A&B\\C&D\end{smallmatrix}\right)\) represents a point via the standard basepoint convention \(Z=S\cdot(iI_g)=(AiI_g+B)(CiI_g+D)^{-1}\), followed by the quotient. The source does not state a representation convention explicitly.

\emph{Problem 4.10 (algorithmic Schottky problem).} Give an algorithm in the sense of complexity theory over \(\R\), with input a \(2g\times2g\) symplectic matrix representing a principally polarized abelian variety, that decides whether the represented point lies in \(\mathcal J_g\). The algorithm must handle the intended exact real inputs; a merely formal algebraic characterization or an unspecified numerical approximation is not the same specification.

A proposed variant fixes \(\eps=\eps(g)>0\) and asks whether the point is within \(\eps\) of \(\mathcal J_g\) for the locally symmetric metric on \(\mathcal A_g\). Writing
\[
 d(x,\mathcal J_g)=\inf_{y\in\mathcal J_g}d(x,y),
\]
this requires deciding the corresponding distance-threshold predicate; the source does not choose strict versus non-strict inequality at \(\eps\). It does not require a polynomial complexity bound or specify whether genus is fixed or part of the input.

The source credits S. Grushevsky [Gr] with making the KP-equations solution effective algebraically, while stressing that this seems different from the requested real-complexity algorithm. Its adjacent discussion studies the distant geometry of the Schottky locus using \(\operatorname{Cone}(\mathcal A_g)\), the tangent cone at infinity defined in its missing equation (10). It credits Hattori [Hat] with identifying that cone with an open cone on a quotient of the Tits boundary and isometrically with a Weyl chamber, a Euclidean sector of dimension \(g\). This is contextual geometry, not a supplied decision algorithm.
\subsection{Short English statement}
Decide exact membership in the period locus from a real symplectic matrix in a specified real-computation model. The source also proposes a fixed-distance variant and distinguishes algebraically effective KP criteria from the desired algorithmic formulation.
\subsection{Sources}
\begin{itemize}
\item Local UnsolvedMath registry, permanent ID 11000047, source code catalogue entry, Problem 4.10, ``Algorithmic Schottky problem,'' PDF page 42. No individual proposer is supplied.
\item Farb, \emph{Problems on Mapping Class Groups and Related Topics}. Supplied URL: \url{https://www.math.uchicago.edu/~farb/papers/mcgbook.pdf}; recorded access 2026-07-29; extraction pdf-font-structure. The edited volume contains 329 explicitly typeset headings across contributed chapters, with repeated numbering disambiguated by page.
\item S. Grushevsky [Gr] is cited for algebraically effective KP equations, and Hattori [Hat] for the asymptotic cone. Their full bibliography and the source equation (10) are absent. The summary additionally mentions Schottky--Jung relations and the Novikov--Shiota KP characterization without full references.
\item Supplied status open and OPEN-TRIAGE. The summary describes partial/classical characterizations and says an algorithmic version is not settled without a precise formulation. This is a supplied assessment. Current primary-literature status and redistribution terms for the publicly accessible author/publisher copy remain NEEDS\_REVIEW.
\end{itemize}
\end{document}
